% Part:sets-functions-relations
% Chapter: size-of-sets
% Section: introduction

\documentclass[../../../include/open-logic-section]{subfiles}

\begin{document}

\olfileid{sfr}{siz}{int}

\olsection{Inleiding}

Toen Georg Cantor in de jaren 1870 de verzamelingenleer ontwikkelde, wilde hij
onder meer het idee van een oneindige verzameling aanvaardbaar maken---een
actuele oneindigheid, zoals de middeleeuwers zouden zeggen. Een belangrijk
onderdeel daarvan was zijn behandeling van de \emph{grootte} van verschillende
verzamelingen. Als $a$, $b$ en $c$ alle verschillend zijn, is de verzameling
$\{a, b, c\}$ intuïtief \emph{groter} dan $\{a, b\}$. Maar hoe zit dat bij
oneindige verzamelingen? Zijn die allemaal even groot? Dat blijkt niet zo te
zijn.

Het eerste belangrijke begrip is dat van een opsomming. Iedere eindige
verzameling kan worden opgesomd door al haar !!{element}s in een lijst te
zetten. Ook van sommige oneindige verzamelingen kunnen we alle !!{element}s
opsommen, mits de lijst zelf oneindig mag zijn. Zulke verzamelingen heten
!!{enumerable}. Cantors verrassende resultaat, dat we aan het einde van dit
hoofdstuk volledig zullen begrijpen, luidt dat sommige oneindige verzamelingen
niet !!{enumerable} zijn.
\end{document}
